% TOP_PROBLEM: {"id": 3226, "title": "Reed's omega, delta, and chi conjecture", "category_id": 3, "category": {"id": 3, "name": "graph_theory", "display_name": "Graph Theory", "description": "Problems involving graphs, networks, and their properties.", "slug": "graph-theory", "order_index": 3, "created_at": "2026-07-31T15:26:25.671Z"}, "status": "open", "problem_number": "OPG-335", "set_id": 12}
% FIRST_POSTED: 2026-10-01
% ATTEMPT_STATUS: unresolved
% MODEL: GPT 6 Astra Ultra
% UnsolvedMath ID 3226; source catalogue code OPG-335
% !TeX program = lualatex
% Research attempt dated 2026-10-01; canonical statement preserved.
\documentclass[11pt,a4paper]{article}
\usepackage[margin=25mm]{geometry}
\usepackage{fontspec}
\setmainfont{lmroman10-regular.otf}[BoldFont=lmroman10-bold.otf,
 ItalicFont=lmroman10-italic.otf,BoldItalicFont=lmroman10-bolditalic.otf]
\usepackage{amsmath,amssymb,mathtools}
\usepackage{unicode-math}
\setmathfont{latinmodern-math.otf}
\AtBeginDocument{\let\setminus\smallsetminus}
\usepackage{xurl,hyperref}
\hypersetup{colorlinks=true,urlcolor=blue}
\setcounter{secnumdepth}{0}
\setlength{\parindent}{0pt}
\setlength{\parskip}{5pt}
\setlength{\emergencystretch}{3em}
\begin{document}
\section*{Reed's omega, delta, and chi conjecture}\label{3226}
{\small UnsolvedMath ID 3226 · OPG-335 · Graph Theory · OpenGarden}
\subsection{Definitions and mathematical statement}
Let $G=(V,E)$ be a nonempty
finite simple undirected graph.
The degree $\deg_G(v)$ of
a vertex counts its neighbours,
and the maximum degree is
$\Delta(G)=\max_{v\in V}\deg_G(v)$.
A clique is a subset of
pairwise adjacent vertices;
write $\omega(G)$ for the
largest clique size. A proper
$q$-colouring assigns one
of $q$ colours to each
vertex so adjacent vertices
receive different colours.
The chromatic number $\chi(G)$
is the least possible $q$.

Reed's conjecture asserts
that every such graph satisfies
\[
 \chi(G)\leq
 \left\lceil\frac{\Delta(G)+1+\omega(G)}2\right\rceil.
\]
The ceiling is applied to
the whole average and means
the least integer no smaller
than that real number.
All three parameters refer
to the same graph.

Every clique requires distinct
colours, so $\omega(G)\leq\chi(G)$.
Greedy colouring gives
$\chi(G)\leq\Delta(G)+1$.
The conjecture says the
rounded-up average of these
two extremal quantities is
itself a universal upper
bound, imposing a stronger
restriction when the maximum
clique is substantially smaller
than the maximum degree.

No connectedness, planarity,
or regularity is assumed.
The conclusion concerns ordinary
vertex colouring, with a
single shared palette; it
does not assert the corresponding
list-colouring inequality.
For an edgeless nonempty graph
the bound reads $1\leq1$.
The empty graph can be
included separately using
chromatic and clique numbers
zero and maximum degree zero.
\subsection{Short English statement}
Is every graph colourable with the rounded-up average of its clique number and one more than its maximum degree?
\subsection{Research attempt}
\paragraph{Scope and method.}
This research attempt by GPT-6 Astra Ultra is dated 1 October 2026.
It preserves the catalogue statement and distinguishes elementary proofs
below from cited prior work. No novelty claim or formal proof-assistant
verification is made. The final paragraph states the precise remaining scope.
\paragraph{Literature and extremal test family.}
Fouquet and Vanherpe [R1] prove Reed's bound for several classes of
expansions of odd cycles, including expansions of a five-cycle.
We give a self-contained certificate for the equal-clique five-cycle
family. It attains the proposed bound at every size and makes the role
of the ceiling explicit. We then derive constraints on a possible
smallest counterexample rather than extrapolate from that family.

\paragraph{Exact graph parameters.}
Let $G_t=C_5[K_t]$, with $t\ge1$: replace each cycle vertex by a clique
$V_i$ of order $t$ and join two such cliques completely precisely when
their indices are consecutive modulo five. There are $5t$ vertices.
Every vertex has $t-1$ neighbours in its own clique and $t$ in each of
two neighbouring cliques, giving
\[
          \Delta(G_t)=3t-1,\qquad \omega(G_t)=2t.
\]
The clique bound is exact because three cycle indices never form a
triangle. An independent set contains at most one vertex per clique
and at most two clique indices. Hence $\alpha(G_t)=2$, and any colouring
needs at least $\lceil5t/2\rceil$ colours.

\paragraph{A matching upper-bound certificate.}
If $t$ is even, for each of the five unordered nonadjacent index pairs
$\{i,i+2\}$ create $t/2$ new colours. Each index belongs to two pairs,
so its clique meets exactly $t$ of these colours. Assign them bijectively
to its $t$ vertices. Every colour is used on exactly two vertices in
nonadjacent cliques, and no colour repeats inside a clique. This is a
proper colouring with $5t/2$ colours.

If $t$ is odd, first colour $t-1$ vertices of each clique by that even
construction, using $5(t-1)/2$ colours. The five remaining vertices induce
$C_5$, which can be coloured using three fresh colours, in cyclic order
$a,b,a,b,c$. The total is $(5t+1)/2=\lceil5t/2\rceil$. Thus
\[
 \chi(G_t)=\left\lceil\frac{5t}{2}\right\rceil
 =\left\lceil\frac{\Delta(G_t)+1+\omega(G_t)}2\right\rceil.
\]
This supplies infinitely many equality cases. For odd $t$ the unrounded
right side is a half-integer, so dropping the ceiling would be false.
The construction is a certificate for every $t$, not just a computation
of small chromatic numbers.

\paragraph{Restrictions on a minimal counterexample.}
Write $R(G)=\lceil(\Delta(G)+1+\omega(G))/2\rceil$. If every connected
component satisfies the bound, their disjoint union does too: its
chromatic number is the maximum of the component chromatic numbers,
and each component's $R$ is at most $R(G)$. A smallest counterexample
would therefore be connected.

It would also be vertex-critical. If deleting a vertex left its chromatic
number unchanged, the smaller graph would satisfy the conjecture by
minimality, and its maximum degree and clique number could not exceed
those of the original graph, a contradiction. Put $h=\chi(G)$ for such
a hypothetical graph. Every vertex has degree at least $h-1$: otherwise
an $(h-1)$-colouring after deleting it could be extended by an unused
colour. Since $h>R(G)$, it follows that $\delta(G)\ge R(G)$. These are
necessary structural conditions, not a construction of a counterexample.

\paragraph{Why the equality-family method stops.}
The colouring of $G_t$ reduces to pairing vertices across five known
cliques. An arbitrary graph need not have that partition, and its
complement may not have enough disjoint nonedges to pair almost all
vertices. Even when large independent sets exist, deleting them changes
maximum degree and clique number by unequal amounts, so an induction
must control the rounded bound carefully. The equality calculation
leaves no universal slack that could absorb an uncontrolled extra colour.

\paragraph{Verification and final status.}
The root batch script constructs the colour classes for $1\le t\le20$,
checks every edge against them, and verifies the parameter formulas.
The independence argument gives the matching lower bound at every size.
\textbf{UNRESOLVED.} The conjecture is proved here for this equality family,
and any minimal counterexample must satisfy the stated criticality and
degree conditions. Neither step bounds the chromatic number of all graphs.

\subsection{Sources}
\begin{itemize}
\item[S1] ULAM.AI and the UnsolvedMath Contributors, supplied problems.json, record 3226 (OPG-335), OpenGarden collection. Catalogue identity and source transcription; CC BY 4.0.
\url{https://huggingface.co/datasets/ulamai/UnsolvedMath}.
\item[S2] Open Problem Garden, Reed's omega, delta, and chi conjecture. Original problem and discussion; the mathematical formulation below is expanded from the supplied source record.
\url{http://www.openproblemgarden.org/op/reeds_omega_delta_and_chi_conjecture}.
\item[R1] Jean-Luc Fouquet and Jean-Marie Vanherpe,
\emph{Reed's Conjecture on hole expansions}, arXiv:1205.0731.
Primary abstract checked 1 October 2026.
\url{https://arxiv.org/abs/1205.0731}.
\end{itemize}
\end{document}
